\documentclass[11pt,a4paper]{article}
\usepackage[margin=27mm]{geometry}
\usepackage{amsmath,amssymb,amsthm,hyperref,booktabs}
\newtheorem{lemma}{Lemma}\newtheorem{theorem}[lemma]{Theorem}
\newcommand{\Sz}{\operatorname{Sz}}\newcommand{\F}{\mathbb F}
\newcommand{\gen}[1]{\langle#1\rangle}
\title{Realizations and obstructions for the displayed patterns}
\author{Human Compression project}\date{}
\begin{document}\maketitle
\begin{abstract}
We assemble the previously reviewed proofs concerning the individual
patterns in the rank-three Suzuki article. Each of its three prime-case
patterns has realizations with the prescribed intersections and RWPRI
property. Among its ten composite-case patterns, four have impossible
intersections, five have RWPRI realizations, and one has the displayed
intersections but fails RWPRI. These conclusions concern the listed
patterns; the source's claim that its list is exhaustive is false and
is not a premise of the arguments below.
\end{abstract}

\section{Coordinates and the two rigidity arguments}
Throughout $q=2^n$ with $n$ odd. We use the classical Suzuki matrix
model, its unique Bruhat coordinates, the usual subfield conjugacy
classes, involution centralizers, cyclic torus centralizers and their
normalizers. Prime relative degree makes an inclusion of simple Suzuki
subfield groups maximal. These are the attributed structural premises
of the original reviewed arguments. The source is D.~Leemans,
\emph{The rank 3 geometries of the simple Suzuki groups $\Sz(q)$},
Note di Matematica 19(1) (1999), 43--63.

Write $S_d=\Sz(2^d)$ for the standard subgroup obtained by restricting
matrix parameters to $\F_{2^d}$. This is initially a field-exponent
index. Unique coordinates give
\[
 S_a\cap S_b=S_{\gcd(a,b)}.
\]
In particular $S_1=\Sz(2)\cong C_5\rtimes C_4$ has order 20; it is
not a simple Suzuki group. Let $H_m$ be the unique subgroup of order
$m$ in the standard split torus, let $w$ invert that torus, and put
$D_d=H_{2^d-1}\rtimes\gen w$ for $d>1$ dividing $n$.
The split-normalizer uniqueness argument, also used in the accepted
normalizer and subfield notes, gives
\[
 N_G(D_d)=D_d,
\]
and exactly one standard-size subfield group of any prescribed degree
containing $d$ contains this embedded $D_d$.

For clarity, the argument for this uniqueness is as follows. A normalizer
of $D_d$ normalizes its characteristic odd subgroup and hence lies in
the full split normalizer. A torus element $h$ normalizes $D_d$ exactly
when $h^2\in H_{2^d-1}$. The quotient torus has odd order, so this is
equivalent to $h\in H_{2^d-1}$. Within a subfield group, every copy of
$D_d$ is conjugate to the standard one: its odd rotation group lies in
a split torus, and squaring in that odd torus aligns the reflection.
Pulling an inclusion back into the standard subfield, then making this
internal conjugation, leaves a conjugator in $N_G(D_d)$. It therefore
does not change the subfield group. This proves the asserted uniqueness.

\begin{lemma}\label{split}
Suppose conjugates $S,T$ of $S_a,S_b$ have full split dihedral
subgroups lying in a common split dihedral subgroup of $G$, and these
two dihedrals share an involution. Then $S,T$ are simultaneously conjugate
to $S_a,S_b$. In particular $S\cap T\cong\Sz(2^{\gcd(a,b)})$.
\end{lemma}
\begin{proof}
Conjugate the common ambient split normalizer to the standard one,
and then conjugate its shared reflection to $w$ using the torus.
The two rotation subgroups are the unique cyclic subgroups of their
orders. The preceding uniqueness argument now forces both subfield
groups to be standard.
\end{proof}

The involution version is equally useful. Fix $z=u(0,1)$ in the
standard centre $Z_q$ of a Sylow two-subgroup $U_q$. Every conjugate
of $S_d$ containing $z$ is $S_d^u$ for some $u\in C_G(z)=U_q$,
because involutions form one class inside $S_d$. Its involutions
commuting with $z$, together with the identity, are precisely
$Z_{2^d}$: conjugation by $U_q$ fixes the entire $Z_q$ pointwise.
Thus if conjugates of $S_a,S_b$ share $z$, their common involutions
commuting with $z$, together with the identity, form
$Z_{2^{\gcd(a,b)}}$. If $a,b$ are coprime, their intersection cannot
contain a Klein four group.

\section{The four impossible composite patterns}
In Theorem~5.5, patterns 1, 2, 3 require distinct odd prime degrees $a,b$,
subgroups $S\cong\Sz(2^a)$, $T\cong\Sz(2^b)$, and a split dihedral
group containing their full split normalizers with a common reflection.
They prescribe $S\cap T$ to have type $D_{10}$, $C_2^2$, or $C_4$.
Lemma~\ref{split} instead forces $S\cap T=\Sz(2)$, of order 20.
Hence all three displayed intersections are impossible. This argument
does not rely on either Mersenne number being prime.

Pattern 5 requires the same coprime field degrees and an intersection
$C_4\times C_2$. That group contains a Klein four group, contradicting
the preceding involution argument. This also independently excludes
pattern 2. These are obstructions for every permitted relative position,
not merely failures of a preferred coordinate construction.

\section{The three patterns when \texorpdfstring{$q-1$}{q-1} is prime}
Here $q-1$ is prime and $q\ge8$. Put $z(x)=u(0,x)$, and
$\theta=2^{(n+1)/2}$. We use the classical facts that every nonidentity
odd-order element belongs to a cyclic torus of order
$q-1,\alpha_q=q+\sqrt{2q}+1$, or $\beta_q=q-\sqrt{2q}+1$,
with normalizer quotients of orders two, four, four. The nonidentity
elements in a Sylow two-subgroup have trace zero.

\begin{lemma}\label{trace}
The $q-1$ matrices $wz(x)$, $x\ne0$, represent all nonidentity
odd-order conjugacy classes, one each.
\end{lemma}
\begin{proof}
Their traces are the distinct nonzero values $x^\theta$. An element
of even order with nontrivial odd part would centralize that odd part,
whose centralizer is an odd cyclic torus; this is impossible. Thus an
even-order element is a two-element, and its trace is zero. Nonzero
trace therefore forces odd order. The displayed matrices belong to
distinct odd-order classes. Torus centralizers and normalizers
give the total number of these classes as
\[
 \frac{q-2}{2}+\frac{\alpha_q-1}{4}+\frac{\beta_q-1}{4}=q-1.
\]
Indeed the normalizer quotients act freely on nonidentity torus elements,
and conjugacy inside a torus is controlled by its normalizer. This proves
exhaustion. In particular every prime dividing $(q-1)(q^2+1)$ occurs
as $|wz(x)|$, and $(q-2)/2$ choices give order $q-1$.
\end{proof}

For the first source pattern, choose involutions $a,b$ with $|ab|=q-1$.
Their unique ovoid fixed points $P,Q$ are distinct; otherwise they would
lie in the same elementary abelian involution centre and commute.
Let $T$ be the pointwise stabilizer of $P,Q$ and let $Z_P,Z_Q$ be
the centres at those points. Set
\[
 G_0=\gen{a,b},\qquad G_1=Z_PT,\qquad G_2=Z_QT.
\]
Then $G_1\cap G_2=T$. No nonidentity rotation of $G_0$ fixes $P$:
its two fixed points are exchanged by $a$, whereas $a$ fixes $P$.
Two reflections fixing $P$ would commute, impossible in an odd dihedral
group. Hence $G_0\cap G_1=\gen a$; similarly
$G_0\cap G_2=\gen b$, and the triple intersection is trivial.

For the second pattern, choose $x$ with $|wz(x)|=p$, where $p$ is
the prescribed odd prime, and take
\[
 G_0=H_{q-1}\rtimes\gen w,\qquad G_1=Z_qH_{q-1},\qquad
 G_2=\gen{w,z(x)}.
\]
The intersections are $H_{q-1},\gen w,\gen{z(x)}$, and the triple
is trivial. To see that the last two intersections are no larger,
a subgroup of $D_{2p}$ containing a reflection is either that subgroup
of order two or the whole dihedral group. The latter would put
$z(x)$ in the split normalizer or $w$ in the point stabilizer,
respectively, both impossible.

For the third pattern, the parabolics are generated by the pairs of
three involutions. The two prescribed half-orders cannot both be two, by the necessary
patterns proved in the preceding note. If one is two, let the other
$p'$ be odd, choose $t,x$
with $|wz(t)|=q-1$, $|wz(x)|=p'$ and $x\ne t$, and use
$w,z(t),z(x)$. If $p'=q-1$ there are at least three available parameters
by Lemma~\ref{trace}. The orders are $2(q-1),4,2p'$. The two central
involutions commute; an odd dihedral group cannot contain two distinct
commuting involutions. Consequently all pair intersections are the
indicated groups of order two, and the triple intersection is trivial.

For distinct odd $p,p'$, choose $x,y$ with $|wz(x)|=p$,
$|wz(y)|=p'$. Choose $\lambda$ satisfying
$\lambda^{\theta+1}x^\theta=y^\theta$ and put
\[
 a=w,\qquad b=h(\lambda)w,\qquad c=z(x).
\]
The exponent $\theta+1$ is coprime to $q-1$, so $\lambda$ exists
uniquely, and $\lambda\ne1$. Thus $ab=h(\lambda)^{-1}$ has order
$q-1$. The traces give $|ac|=p$, $|bc|=p'$ by Lemma~\ref{trace}.
An intersection properly containing a shared reflection would identify
the corresponding prime dihedral groups, which is impossible from their
different orders or would put $c$ in the split normalizer. Thus the
required intersections again hold.

For equal odd $p=p'$, choose distinct $t,x$ with $|wz(t)|=q-1$,
$|wz(x)|=p$ and take
\[
 a=w,\qquad b=z(t)wz(t),\qquad c=z(x).
\]
Now $ab=(wz(t))^2$ has order $q-1$, and $bc$ is conjugate to $ac$
because $z(t)$ commutes with $c$. Also
$\gen{a,b}=\gen{w,z(t)}$, while $c$ does not belong to that dihedral
group, since it is distinct from and commutes with $z(t)$. The same
maximal-reflection argument gives all intersections, including $p=5$.

In each pattern the full split normalizer is maximal in $G$.
The reflection intersections have prime index in each odd dihedral
parabolic. The torus complement is maximal in $Z_qH_{q-1}$ by its
irreducible scalar action. Finally all nontrivial pair intersections
over the trivial triple have prime order. These facts verify every
subset in the source RWPRI axiom. We have proved existence of the
displayed subgroup patterns, without inferring thinness or flag transitivity
for all their possible embeddings.

\section{The realizable composite patterns}
For this section only, $S_s$ denotes the standard group $\Sz(s)$
with a field-size index. Patterns 4 and 6 are the following one family.
Let $q=s^r$, $s\ge8$, $r$ an odd prime, choose a prime
$\ell\mid(q-1)/(s-1)$ and put $m=(s-1)\ell$. Take
\[
 G_0=S_s,\qquad G_1=H_m\rtimes\gen w,\qquad G_2=Z_qH_m.
\]
The intersections are respectively
$D_{2(s-1)},Z_sH_{s-1},H_m$, with full intersection $H_{s-1}$.
Unique coordinates prove these equalities. The subfield and its full
split normalizer give maximal inclusions, and the cyclic extension
indices are $\ell$. The complement $H_m$ is maximal in $Z_qH_m$:
its scalar field contains $\F_s$, cannot equal $\F_s$ because
$m>s-1$, and has degree dividing the prime $r$. Thus it is $\F_q$,
and the scalar action is irreducible. These facts verify all RWPRI
subsets. The torus action on $Z_q$ uses a power coprime to $q-1$,
which does not change the generated scalar field.

For patterns 7, 8, 9 put $q=t^{ab}$, $s=t^a$, $s'=t^b$, where
$t\ge8$ and $a,b$ are distinct odd primes. For pattern 7 choose distinct
primes $p\mid(s-1)/(t-1)$, $p'\mid(s'-1)/(t-1)$ and put
$m=(t-1)pp'$. Use $S_s,S_{s'}$ and $D=H_m\rtimes\gen w$.
The cyclic gcd identities are
\[
 \gcd(m,s-1)=(t-1)p,\qquad
 \gcd(m,s'-1)=(t-1)p'.
\]
They hold also when a selected prime divides $t-1$: its additional
valuation occurs in exactly its selected relative prime degree.
Equivalently, these are the two-index instances of the scalar-field
arithmetic lemma in the separately reviewed split-family note.
The pair intersections are $S_t,D_{2(t-1)p},D_{2(t-1)p'}$, and the
full intersection is $D_{2(t-1)}$. It is maximal in $S_t$, and has
prime index in each larger dihedral; all subfield inclusions have
prime relative degree. Hence RWPRI holds.

For pattern 8 choose a full nonsplit torus $A_t$ of $S_t$, of order
$m=t\pm\sqrt{2t}+1$, and an element $r$ of order four in its
normalizer. Its ambient centralizer is one cyclic nonsplit torus $A_q$.
Put $A_s=C_{S_s}(A_t)$ and $A_{s'}=C_{S_{s'}}(A_t)$.
Their intersection is $A_t$, so their orders have gcd $m$.
Choose primes $p,p'$ with $mp\mid|A_s|$, $mp'\mid|A_{s'}|$.
They are distinct, and the corresponding subgroup $C_{mpp'}$ exists
in $A_q$. Then
\[
 K=\gen{C_{mpp'},r},\qquad
 K\cap S_s=C_{mp}\rtimes\gen r,\qquad
 K\cap S_{s'}=C_{mp'}\rtimes\gen r,
\]
with full intersection $A_t\rtimes\gen r$. The same prime-index
and maximal-normalizer checks prove RWPRI. The compatible nonsplit
torus sign is determined by $A_t$, rather than chosen independently
in the two subfields.

For pattern 10 return to exponent indices. Choose distinct prime
divisors $a,b,c$ of $n$, with $n/(abc)>1$ as in the printed parameter
condition. The standard subfields
$S_{n/a},S_{n/b},S_{n/c}$ have exactly the displayed pair and triple
intersections by field gcds. All one-step inclusions have prime
relative degree between simple subfield groups, so every RWPRI
condition holds.

\section{The mixed pattern: intersections without maximality}
Pattern 9 prescribes $S_s,S_{s'}$ with common $S_t$ and a third group
of type $E_q\mathbin{\hat\cdot}E_t\rtimes C_{t-1}$.
Its intersections with the two simple groups have types
$E_s\mathbin{\hat\cdot}E_t\rtimes C_{t-1}$ and
$E_{s'}\mathbin{\hat\cdot}E_t\rtimes C_{t-1}$; the full
intersection is the point stabilizer $B_t=U_tH_{t-1}$.
Here the first elementary abelian factor has the literal full order
$q$, as printed; it is not an unspecified smaller subspace.

Conjugate the common $S_t$ to standard position. It contains its
full split normalizer, so the uniqueness argument forces the two
containing simple groups to be the standard $S_s,S_{s'}$.
An internal conjugation in $S_t$ puts their common Borel at $B_t$.
The third parabolic has a normal Sylow two-subgroup $P$ of order $qt$
containing an elementary abelian group of order $q$. In an ambient
Sylow two-subgroup every involution is central, so this elementary
abelian group is its full centre. It consists of the identity and all
involutions of $P$, hence is characteristic in $P$ and normal in the
third parabolic. Its unique fixed ovoid point is therefore preserved.
Containment of $B_t$ forces this point to be the standard point at
infinity. Thus the normal elementary abelian group is $Z_q$, and
containment and orders force the third parabolic and its intersections
to be
\[
 Z_qU_tH_{t-1},\qquad Z_sU_tH_{t-1},\qquad
 Z_{s'}U_tH_{t-1}.
\]
Conversely these standard groups have all the printed intersections,
so the bare subgroup pattern exists.

For any $\F_t$-subspace $V$ with $\F_s<V<\F_q$, the group
$Z_VU_tH_{t-1}$ lies strictly between the first intersection and the
third parabolic. Such a subspace exists: adjoining one vector to
$\F_s$ gives dimension $a+1<ab$. The same argument gives an
intermediate group above the second intersection, using $b+1<ab$.
Hence neither remaining intersection is maximal in this parabolic.
The RWPRI condition for its singleton label fails. The forced
coordinates show that this is an obstruction for every realization
of the printed type, not just the one displayed here.

\section{What these arguments settle}
The original final paragraph asks about actual existence of the
displayed intersections. The retained results are
\[
\begin{array}{l|c|c}
\text{source patterns}&\text{intersections exist}&\text{RWPRI realization}\\\hline
\text{Theorem 5.3, all three}&\text{yes}&\text{yes}\\
\text{Theorem 5.5, }1,2,3,5&\text{no}&\text{no}\\
\text{Theorem 5.5, }4,6,7,8,10&\text{yes}&\text{yes}\\
\text{Theorem 5.5, }9&\text{yes}&\text{no}
\end{array}
\]
These are the previously reviewed phase-five arguments, reassembled
with the now accepted uniqueness and scalar-field interfaces made
explicit. Their old sufficient-pattern proofs did not establish the
source's completeness claim or Theorem 5.4; the present source treatment
uses separate reviewed proofs and counterexamples for those assertions.
No classification of all relative positions or all flag-transitive
geometries is inferred from this table. The saved finite examples
support the constructions but are not substituted for the uniform
proofs, and their acquisitions are not repeated here.
\end{document}
